%%%%%
%
\section{Filtre num�rique �quivalent � un circuit �lectrique RC}

\itemstyle{Op\'erateur~:} {\sf Filtre\_RC}
\index{Filtre\_RC}

%%%%%OBJET%%%%
\itemstyle{Objet~:}
Cet op�rateur applique � un param\`etre un filtrage num�rique �quivalent
� un circuit �lectrique RC. Il est d�fini par sa constante de temps $\tau$\\ 
Les caract�ristiques du filtre sont les suivantes : \\
\begin{itemize}
\item R�ponse � un cr�neau d'amplitude $X_0$ : 
$X_f = X_0 . ( 1 - e^{-\frac{t}{\tau}} )$   
\item Fonction de transfert en transform�e de Laplace :  
$ H(p) = \frac{1/\tau}{p.(p+1/\tau)} $
\item Fonction de transfert en transform\'ee en Z :
$H(Z) = \frac{1 - e^{-Te/\tau} }{ Z - e^{ -Te/\tau} } $ 
\item Fr�quence de coupure � 3dB : $f_c = \frac{1}{2\pi\tau} $ 
\end{itemize}

%%%%%ARGUMENTS%%%%
\bigskip
%%%%%ENTREES%%%%
\begin{tabarg}
\multicolumn{3}{l}{\titretab {Entr\'ees}}\\
\hline
$X$ & Vecteur, {\tt float} & Param�tre d'entr�e \\
$f_e$ & Coeff., {\tt float} & Fr�quence d'�chantillonnage de $X$ (Hz)\\
$\tau$ & Coeff., {\tt float} & Constante de temps du filtre ($s$)\\
\end{tabarg}

\medskip
%%%%%SORTIES%%%%
\begin{tabarg}
\multicolumn{3}{l}{\titretab {Sorties}}\\
\hline
$X_f$ & Vecteur, {\tt float} & Param\`etre filtr\'e (m\^eme unit\'e que $X$)\\
\end{tabarg}


%%%%%FORMULE%%%%
\itemstyle{Formule}

%%Filtre_RC_sdph
\[
X_f[i] = e^{-\frac{1}{f_e.\tau}}.X_f[i-1] +
         \left( 1 - e^{-\frac{1}{f_e.\tau}} \right). X[i-1]
\]



%%%%%AUTEUR%%%%
\itemstyle{Auteur~:}
CNRM/GMEI/TRAMM

%%%%%REFERENCE%%%%
\itemstyle{R\'ef\'erences~:}
